\setlength{\tabcolsep}{2.5pt}
\begin{tabular}{@{}lp{0.86\textwidth}@{}}
\toprule
Filler & Text \\
\midrule
\multicolumn{2}{@{}l}{\slotmarker{S1} sign slot, $a_{4}$--$a_{5}$} \\
\quad conflict & Curtis cited novae observed in spiral nebulae as evidence that those nebulae lay at great distances; that evidence indicates a remove for M31 of about 10,000 light years from the Sun, the figure his analysis of those novae yielded. \\
\quad no link & Curtis cited novae observed in spiral nebulae as evidence that those nebulae lay at great distances. His analysis of novae in spiral nebulae placed M31 about 10,000 light years from the Sun. \\
\quad support & Curtis cited novae observed in spiral nebulae as evidence that those nebulae lay at great distances, and that evidence rules out the claim, sometimes attached to his analysis of novae in spiral nebulae, that it placed M31 about 10,000 light years from the Sun. \\
\addlinespace
\multicolumn{2}{@{}l}{\slotmarker{S2} sign (spare) slot, $a_{7}$--$a_{8}$} \\
\quad conflict & Shapley relied on Adriaan van Maanen's reported measurements of internal rotation in spiral nebulae, and that reliance shows that the rotations van Maanen reported were confirmed by later comparisons of photographic plates. \\
\quad no link & Shapley relied on Adriaan van Maanen's reported measurements of internal rotation in spiral nebulae. The rotations van Maanen reported were confirmed by later comparisons of photographic plates. \\
\quad support & Shapley relied on Adriaan van Maanen's reported measurements of internal rotation in spiral nebulae, and because nothing beyond those original reports of rotation stood behind the figures, the rotations cannot have been confirmed by later comparisons of photographic plates. \\
\addlinespace
\multicolumn{2}{@{}l}{\slotmarker{S3} control slot, $a_{11}$--$a_{12}$} \\
\quad wording 1 & Edwin Hubble identified Cepheid variable stars in M31 using the 100-inch Hooker telescope at Mount Wilson, and that identification led him to a distance for the nebula of roughly 900,000 light years. \\
\quad wording 2 & Because Edwin Hubble was able to pick out Cepheid variable stars in M31 with the 100-inch Hooker telescope at Mount Wilson, a distance of roughly 900,000 light years followed for that nebula. \\
\quad wording 3 & A distance to M31 of roughly 900,000 light years was the direct result of Edwin Hubble's detection of Cepheid variable stars in that nebula with the 100-inch Hooker telescope at Mount Wilson. \\
\bottomrule
\end{tabular}
